\section{Checklist for building a probability model}

The examples in this chapter have different physical forms: coins, dice, cards,
urns, and waiting times. The same modelling discipline appears in all of them.
We do not begin with a formula. We begin by constructing the probability space.

The purpose of this final checklist is to make that construction explicit.
When facing a new probability problem, the following questions should be answered
before any numerical calculation is attempted.

\subsection*{Step 1: What is the experiment?}

Start with the physical or conceptual procedure.

Examples include:
\begin{itemize}
    \item toss a coin once;
    \item toss a coin \(n\) times;
    \item roll two dice;
    \item draw five cards from a deck;
    \item select objects from an urn;
    \item repeat trials until the first success occurs.
\end{itemize}

At this stage the description may still be incomplete. The phrase ``draw five
cards'' does not say whether order is recorded. The phrase ``select from an
urn'' does not say whether replacement occurs. The phrase ``repeat trials'' does
not say what is recorded when the process stops.

\subsection*{Step 2: What is recorded?}

The recording rule determines the mathematical outcome.

Ask:
\begin{itemize}
    \item Is order recorded?
    \item Are labels recorded, or only types such as colour or rank?
    \item Is the whole sequence recorded, or only a numerical summary?
    \item Is the experiment performed with replacement or without replacement?
    \item Does the experiment stop after a fixed number of trials, or at a random time?
\end{itemize}

This step is essential because the same physical procedure can produce different
sample spaces under different recording rules.

\begin{remarkbox}
Most elementary probability errors begin when the recording rule is left
implicit. If it is not clear what is recorded, then it is not clear what one
elementary outcome is.
\end{remarkbox}

\subsection*{Step 3: What is one elementary outcome?}

Once the recording rule is fixed, describe one complete elementary outcome.

Typical forms include:
\begin{itemize}
    \item a single symbol, such as \(H\) or \(T\);
    \item a number, such as a die result \(i\in\{1,
    \ldots,6\}\);
    \item an ordered pair, such as \((i,j)\) for two dice;
    \item a sequence, such as \(HTTH\);
    \item a set, such as a five-card hand;
    \item a count, such as the number of heads;
    \item a stopping time, such as the trial number of the first success.
\end{itemize}

The elementary outcome must be complete relative to the recording rule. If the
rule records the full sequence of coin tosses, then the number of heads is not a
complete elementary outcome. If the rule records only the number of heads, then
the full sequence contains more information than the model keeps.

\subsection*{Step 4: What is the sample space \(\Omega\)?}

The sample space is the set of all elementary outcomes:
\[
\Omega=\{\text{all possible elementary outcomes under the chosen recording rule}\}.
\]

Examples:
\[
\{H,T\},\qquad \{H,T\}^n,\qquad \{1,2,3,4,5,6\},
\]
\[
\{1,
\ldots,6\}^2,\qquad \{S\subseteq D: |S|=5\},
\]
\[
\{1,2,3,
\ldots\}.
\]

The sample space may be finite, such as \(\{H,T\}^n\), or countably infinite,
such as the waiting-time space \(\{1,2,3,\ldots\}\).

\subsection*{Step 5: What events are being considered?}

An event is a subset of \(\Omega\). Translate each verbal statement into a set of
elementary outcomes.

Examples:
\[
\text{exactly }k\text{ heads}
\]
becomes
\[
\{\omega\in\{H,T\}^n: \omega\text{ contains exactly }k\text{ heads}\}.
\]
The statement
\[
\text{the sum of two dice is }7
\]
becomes
\[
\{(1,6),(2,5),(3,4),(4,3),(5,2),(6,1)\}.
\]
The statement
\[
\text{a five-card hand contains at least one ace}
\]
becomes a set of five-card subsets of the deck.

\begin{remarkbox}
Do not count before the event has been written clearly. Counting favourable
outcomes requires knowing exactly which outcomes are favourable.
\end{remarkbox}

\subsection*{Step 6: Are elementary outcomes equally likely?}

This is the question that determines whether the classical formula applies.

If \(\Omega\) is finite and all elementary outcomes are equally likely, then
\[
P(A)=\frac{|A|}{|\Omega|}.
\]
This applies, for example, to:
\begin{itemize}
    \item a fair coin tossed \(n\) times, with full sequence recorded;
    \item two fair dice, with ordered pair recorded;
    \item a well-shuffled five-card hand, with the hand recorded as a set.
\end{itemize}

If elementary outcomes are not equally likely, then the formula
\[
P(A)=\frac{|A|}{|\Omega|}
\]
must not be used. Instead, add the probabilities of the elementary outcomes in
\(A\):
\[
P(A)=\sum_{\omega\in A}P(\omega).
\]

Examples of non-uniform models include:
\begin{itemize}
    \item a biased coin;
    \item a loaded die;
    \item sums of two fair dice when the sample space is \(\{2,3,\ldots,12\}\);
    \item colours recorded from an urn with unequal numbers of different types;
    \item waiting time until first success.
\end{itemize}

\subsection*{Step 7: Count or sum carefully}

In a uniform finite model, the main task is often to count:
\[
|A| \quad\text{and}\quad |\Omega|.
\]
The counting method must match the structure of the sample space.

Use sequences when order matters. Use combinations when only membership matters.
Use product spaces for repeated independent stages. Use complement events when
the complement is easier to count.

In a non-uniform model, the main task is not only to count outcomes, but to
assign and sum their probabilities. For example, for a biased coin,
\[
P(\text{exactly }k\text{ heads in }n\text{ tosses})
=
\binom nk p^k(1-p)^{n-k}.
\]
The binomial coefficient counts the sequences, while the factor
\(p^k(1-p)^{n-k}\) gives the probability of each such sequence.

\subsection*{Step 8: Interpret the answer}

A final numerical answer is not the end of the solution. It should be connected
back to the original model.

For example,
\[
\frac6{36}=\frac16
\]
for the event ``sum seven on two fair dice'' means that six of the thirty-six
ordered pairs have sum seven. It does not mean that the value seven is one of
six equally likely sums.

Similarly,
\[
\frac35
\]
for drawing type \(R\) from an urn with three objects of type \(R\) and two of
type \(B\) means that three of the five equally selectable labelled objects
produce the recorded type \(R\). It does not mean that the colour sample space is
uniform.

\subsection*{A compact modelling checklist}

Before calculating, write down the following:

\begin{enumerate}
    \item What is the experiment?
    \item What exactly is recorded?
    \item What is one elementary outcome?
    \item What is the sample space \(\Omega\)?
    \item What event \(A\subseteq\Omega\) is being studied?
    \item Are elementary outcomes equally likely?
    \item If yes, what are \(|A|\) and \(|\Omega|\)?
    \item If no, what are the elementary probabilities \(P(\omega)\)?
    \item How is \(P(A)\) computed from those outcomes?
    \item What does the answer mean in the original situation?
\end{enumerate}

\begin{theorembox}
The basic rule of this chapter is:
\[
\text{construct the probability model first, calculate second.}
\]
A correct probability calculation is a calculation inside a clearly described
sample space with a clearly described probability assignment.
\end{theorembox}

The examples in this chapter should be read as models of reasoning, not as a
list of formulas. The formulas are useful because they preserve the structure of
the model. When the model changes, the formula may change as well.
